\section{Υποστήριξη \ENG{Learning Standards} και Επεκτασιμότητα}
\label{sec:standards_extensibility}

Η υποστήριξη πολλαπλών προτύπων στο παρόν σύστημα δεν υλοποιείται μέσω πολλαπλών ανεξάρτητων \ENG{runtime cores}, αλλά μέσω ενός κοινού πυρήνα που αναθέτει την ερμηνεία του \ENG{payload} σε \ENG{standard-specific adapters}. Ο \texttt{\ENG{StandardAdapterRegistry}} συλλέγει όλες τις υλοποιήσεις του \texttt{\ENG{StandardAdapter}} και τις εκθέτει με κλειδί το \ENG{LearningStandard}. Έτσι, ο \texttt{\ENG{RuntimeOrchestrator}} δεν περιέχει \ENG{if/else} αλυσίδες για κάθε πιθανό πρότυπο, αλλά ζητά τον κατάλληλο \ENG{adapter} με βάση το \ENG{standard} του \ENG{launch}.

Οι \texttt{\ENG{Scorm12Adapter}} και \texttt{\ENG{Scorm2004Adapter}} δείχνουν με σαφήνεια το όριο της τρέχουσας υλοποίησης. Και οι δύο \ENG{adapters} συγχωνεύουν την ακατέργαστη κατάσταση εκτέλεσης της συνεδρίας και έπειτα ενημερώνουν ένα αντικείμενο \texttt{\ENG{NormalizedProgress}}. Στο \ENG{SCORM 1.2} χαρτογραφούνται κυρίως τα \texttt{\ENG{lesson\_status}}, \texttt{\ENG{score.raw}}, \texttt{\ENG{session\_time}}, \texttt{\ENG{total\_time}}, \texttt{\ENG{lesson\_location}} και \texttt{\ENG{suspend\_data}}. Στο \ENG{SCORM 2004} προστίθενται τα \texttt{\ENG{completion\_status}}, \texttt{\ENG{success\_status}}, \texttt{\ENG{score.scaled}}, \texttt{\ENG{score.min}}, \texttt{\ENG{score.max}}, \texttt{\ENG{location}} και \texttt{\ENG{suspend\_data}}. Η υλοποίηση αυτή αρκεί για συγκεντρωτική πρόοδο, αλλά δεν τεκμηριώνει πλήρη κάλυψη του συνολικού \ENG{CMI} μοντέλου.

Αυτό το σημείο είναι θεωρητικά κρίσιμο, επειδή η προδιαγραφή του \ENG{SCORM 2004} έχει ήδη παρουσιαστεί ως πολύ ευρύτερη στα προηγούμενα κεφάλαια. Ο παρών \ENG{engine} ανιχνεύει ότι ένα πακέτο ανήκει στο \ENG{SCORM 2004}, εκθέτει τον κατάλληλο \ENG{API} και ερμηνεύει ένα υποσύνολο πεδίων προόδου, αλλά δεν προκύπτει από τον κώδικα πλήρης μηχανή \ENG{Sequencing and Navigation}, ούτε αποθήκευση \ENG{objectives}, \ENG{activity tree} ή \texttt{\ENG{adl.nav}} αιτημάτων.

Παρά τον περιορισμό αυτό, η επιλογή του \ENG{adapter pattern} έχει πρακτική σημασία. Εάν στο μέλλον προστεθεί επιπλέον πρότυπο ή βαθύτερη κάλυψη κάποιου υπάρχοντος, η κύρια αλλαγή μπορεί να εντοπιστεί σε νέα υλοποίηση \texttt{\ENG{StandardAdapter}} και, ενδεχομένως, σε πλουσιότερο μοντέλο κανονικοποίησης. Ο \ENG{orchestration} πυρήνας, το \ENG{security model}, οι \ENG{launch} μηχανισμοί και τα όρια μόνιμης αποθήκευσης μπορούν να παραμείνουν σταθερά. Αυτό ακριβώς αποτυπώνει το Σχήμα~\ref{fig:adapter_architecture}: η επεκτασιμότητα αφορά τη μετάφραση της σημασιολογίας εκτέλεσης κάθε προτύπου προς ένα κοινό εσωτερικό σχήμα και όχι την ύπαρξη πολλαπλών, πλήρως ανεξάρτητων υποσυστημάτων.

\begin{figure}[H]
  \centering
\begin{chapterfigurebox}
  \begin{tikzpicture}[
      font=\footnotesize,
      node distance=1.8cm and 2.2cm,
      box/.style={draw, rounded corners, fill=gray!5, align=center, text width=3.3cm, minimum height=1.05cm},
      emph/.style={draw, rounded corners, fill=gray!12, align=center, text width=3.6cm, minimum height=1.05cm, font=\footnotesize\bfseries},
      arrow/.style={-Latex, thick},
      boundary/.style={draw, dashed, rounded corners, inner sep=11pt}
    ]
    \node[emph] (runtimeo) {\texttt{\ENG{RuntimeOrchestrator}}};
    \node[box, below=1.8cm of runtimeo] (registry) {\texttt{\ENG{StandardAdapterRegistry}}};

    \node[box, below left=1.9cm and 2.3cm of registry] (s12) {\texttt{\ENG{Scorm12Adapter}}};
    \node[box, below=1.9cm of registry] (s2004) {\texttt{\ENG{Scorm2004Adapter}}};
    \node[box, below right=1.9cm and 2.3cm of registry] (future) {\ENG{Future}\\\texttt{\ENG{StandardAdapter}}};

    \node[emph, below=2.0cm of s2004] (norm) {\texttt{\ENG{NormalizedProgress}}};

    \draw[arrow] (runtimeo) -- (registry);
    \draw[arrow] (registry) -- (s12);
    \draw[arrow] (registry) -- (s2004);
    \draw[arrow, dashed] (registry) -- (future);
    \draw[arrow] (s12) -- (norm);
    \draw[arrow] (s2004) -- (norm);
    \draw[arrow, dashed] (future) -- (norm);

    \node[boundary, fit=(s12)(s2004)(future)] (adapterboundary) {};
    \node[font=\footnotesize, align=center, text width=5.6cm, above=0.14cm of adapterboundary] {\ENG{Pluggable standard interpretation layer}};
  \end{tikzpicture}
\end{chapterfigurebox}
  \caption{Επεκτασιμότητα του \ENG{runtime} μέσω \ENG{registry} και \ENG{standard-specific adapters} που τροφοδοτούν κοινό μοντέλο κανονικοποιημένης προόδου.}
  \label{fig:adapter_architecture}
\end{figure}

